\section*{Capítulo \ref{ch:b3:quantum-statistics} --- Estadística cuántica}

\begin{solution}{exo:b3:quantum-statistics:1}
(a) $\qty{7.1}{eV}$, $\qty{8.2e4}{K}$. (b) $\qty{3.2}{eV}$,
$\qty{3.7e4}{K}$. (c) Con $m = 3u$: $E_{\text{F}} \approx
\qty{4e-4}{eV}$, $T_{\text{F}} \approx \qty{5}{K}$ (valor de gas
ideal; las interacciones lo bajan). (d) Los electrones de los metales:
degenerados a toda temperatura terrestre; el $^3$He: solo por debajo de unos
pocos kelvin; $T = \qty{300}{K}$ nunca roza la escala de Fermi de un metal.
\end{solution}

\begin{solution}{exo:b3:quantum-statistics:2}
(a) $\int_0^{E_{\text{F}}}\epsilon\,g\,\dd\epsilon/\int g\,\dd
\epsilon = \tfrac35 E_{\text{F}}$ a partir de $g \propto \sqrt\epsilon$.
(b) $v_{\text{F}} = \sqrt{2E_{\text{F}}/m_{\text{e}}} =
\qty{1.6e6}{m/s}$; el valor cuadrático medio $\sqrt{3/5}\,v_{\text{F}} \approx
\qty{1.2e6}{m/s}$. (c) $\tfrac32 k_{\text{B}}T = \tfrac35
E_{\text{F}}$: $T \approx \qty{3.3e4}{K}$. (d) No hay adónde
ir: todo estado inferior está ocupado, así que no puede emitirse ningún fotón
ni frenarlos ninguna colisión --- velocidad sin calor, por exclusión.
\end{solution}

\begin{solution}{exo:b3:quantum-statistics:3}
(a) $P = 0.4 \times \num{8.5e28} \times \qty{7.1}{eV} =
\qty{3.9e10}{Pa} \approx \num{4e5}$ atmósferas. (b) La atracción de
Coulomb de la red de iones: el metal es la tregua entre las
dos (\cref{exo:b3:quantum-statistics:11}). (c) Comprimir sube
$E_{\text{F}}$ de forma abrupta (como $n^{2/3}$): los metales se resisten. (d) $K \sim
\tfrac53 P \approx \qty{65}{GPa}$ frente a los
\qty{140}{GPa} medidos: la mitad de la rigidez del cobre es exclusión pura.
\end{solution}

\begin{solution}{exo:b3:quantum-statistics:4}
(a) $T_{\text{c}} \approx \qty{0.4}{\micro K}$ --- el condensado del
JILA apareció a \qty{170}{nK}: la misma física, con una densidad de pico
algo menor. (b) $T_{\text{c}} \propto 1/m$: $\times87 \approx
\qty{35}{\micro K}$. (c) A densidad alta y con ese frío, todos los gases
salvo el hidrógeno polarizado en espín \emph{solidifican}; las colisiones a
tres cuerpos forman moléculas. El truco es un gas diluido y metastable,
lo bastante frío para que $\lambda_T$ compense la dilución. (d)
$\lambda_{T_{\text{c}}} = h/\sqrt{2\pi mk_{\text{B}}T_{\text{c}}}
\approx \qty{0.3}{\micro m}$: $n\lambda^3 = \num{e20} \times
\num{2.7e-20} \approx 2.7$, como exige la regla.
\end{solution}

\begin{solution}{exo:b3:quantum-statistics:5}
(a) $\gamma T = \beta T^3$ en $T^2 = \gamma/\beta$: para el cobre,
$T \approx \qty{3}{K}$. (b) $C_{\text{el}}/C \approx
(\pi^2/2)(300/\num{8.2e4})/3 \approx 0.6\%$. (c) A \qty{0.1}{K},
el término de red se ha desplomado como $T^3$: los electrones aportan
$\sim10^3$ veces más --- la calorimetría se vuelve espectroscopia
electrónica. (d) $C/T = \gamma + \beta T^2$ es una recta en
$T^2$: la ordenada en el origen $\gamma$ mide la \omterm{prop:b3:schrodinger-three-dimensions:counting}{densidad de estados} en el
nivel de Fermi, y la pendiente $\beta$ la temperatura de Debye --- dos números
por metal a partir de una sola gráfica.
\end{solution}

\begin{solution}{exo:b3:quantum-statistics:6}
(a) $N_{\text{e}} = M/2m_{\text{n}} = \num{6e56}$; $n =
\num{5.5e35}\,\unit{m^{-3}}$; $E_{\text{F}} \approx
\qty{0.24}{MeV}$. (b) La mitad de la \omterm{def:b3:relativistic-dynamics:momentum-energy}{energía en reposo} del electrón (al borde
de lo relativista) y, aun así, trescientas veces $k_{\text{B}}T$
(\qty{0.9}{keV}): al rojo blanco y cuánticamente fría a la vez. (c)
$P_{\text{deg}} \approx \qty{9e21}{Pa}$ frente a $GM^2/R^4 \sim
\num{e22}\,\unit{Pa}$: el mismo orden --- la estrella es exactamente ese
equilibrio. (d) Con más masa: el lado de la gravedad crece como $M^2$ y la
estrella debe encogerse ($R \propto M^{-1/3}$), lo que empuja $E_{\text{F}}$
hacia el ablandamiento relativista.
\end{solution}

\begin{solution}{exo:b3:quantum-statistics:7}
(a) $\rho \approx \qty{5e17}{kg/m^3}$: unas dos veces la densidad
nuclear --- un núcleo del tamaño de una ciudad. (b) $n_{\text{n}}
\approx \qty{3e44}{m^{-3}}$: $E_{\text{F}} \approx \qty{90}{MeV}$,
una décima parte de $m_{\text{n}}c^2$: la relatividad llama a la puerta. (c) La
\omterm{thm:b3:quantum-statistics:fermi}{energía de Fermi} de los electrones superó el coste de \qty{0.78}{MeV} de
$\text{e} + \text{p} \to \text{n} + \nu$: la captura se volvió
rentable, y la estrella cambió sus electrones y protones por
neutrones. (d) $\qty{5}{cm^3} \times \qty{5e17}{kg/m^3} =
\qty{2.5e12}{kg}$: una cucharadita que pesa más que una cordillera.
\end{solution}

\begin{solution}{exo:b3:quantum-statistics:8}
(a) Enfriar una cavidad simplemente \emph{quita} fotones (las paredes los
absorben): sin ningún número que conservar, $\mu$ queda clavado en cero y
ningún excedente se amontona jamás en el modo fundamental --- el gas se
apaga en vez de condensar. (b) El número de átomos se conserva en toda
escala de tiempo de laboratorio. (c) El colorante absorbe y reemite una y otra
vez los fotones \emph{sin perderlos} en la escala de tiempo de la
termalización: un número de fotones efectivamente conservado (y una
masa efectiva procedente de la cavidad) --- y la condensación aparece en efecto.
(d) Los fonones: creados y destruidos sin trabas, $\mu = 0$: no hay
condensación --- el ``gas de fonones'' de un cristal simplemente se congela.
\end{solution}

\begin{solution}{exo:b3:quantum-statistics:9}
(a) $T_{\text{c}}^{\text{ideal}} \approx \qty{3.1}{K}$ frente al
$T_\lambda = \qty{2.17}{K}$ medido. (b) El helio líquido es denso
y fuertemente interactuante: ``ideal'' descarta las interacciones que
desplazan (y remodelan) la transición. (c) Los átomos de $^3$He son fermiones:
no hay condensación sin apareamiento; las interacciones atractivas ligan
pares de tipo Cooper solo por debajo de \qty{2.6}{mK}, y el líquido apareado
superfluye entonces. (d) Nueve fermiones: el $^6$Li es un fermión (mientras
que el $^7$Li es un bosón) --- el par de isótopos del litio permite a un mismo
laboratorio estudiar un mar de Fermi y un condensado en el mismo horno.
\end{solution}

\begin{solution}{exo:b3:quantum-statistics:10}
(a) Sustitúyase $g \propto \sqrt\epsilon$ y sáquese $\beta$ del integrando.
(b) El valor de la integral da $N_{\text{exc}} =
2.612\,V/\lambda_T^3$; imponer $N_{\text{exc}} = N$ define
$T_{\text{c}}$. (c) Por debajo de $T_{\text{c}}$, $N_{\text{exc}}(T) =
N(T/T_{\text{c}})^{3/2}$; el resto es $N_0$. (d) En dos
dimensiones, $g(\epsilon)$ es constante y $\int\dd\epsilon/(\eu^{
\beta\epsilon} - 1)$ diverge por abajo: los estados excitados
siempre pueden absorberlos a todos --- no hay población macroscópica del
estado fundamental a ninguna $T > 0$.
\end{solution}

\begin{solution}{exo:b3:quantum-statistics:11}
(a) $\dd\epsilon/\dd n = \tfrac23 An^{-1/3} - \tfrac13 Bn^{-2/3} =
0$ en $n^{1/3} = B/2A$: un mínimo genuino. (b) Con $A \sim
\hbar^2/m_{\text{e}}$ y $B \sim e^2/4\pi\varepsilon_0$, el
espaciado de equilibrio $n^{-1/3} \sim \hbar^2 4\pi\varepsilon_0/
m_{\text{e}}e^2 = a_0$: los metales son densos porque lo dice el \omterm{thm:b3:hydrogen-atom:levels}{radio de
Bohr}. (c) $\epsilon_{\min} \sim -B^2/A \sim -E_{\text{I}}$
--- cohesión de un electronvoltio, tal como se mide. (d) La \omterm{thm:b3:quantum-statistics:fermi}{presión de degeneración}
es la mitad repulsiva de la tregua: sin ella, el término de Coulomb
aplastaría la red hasta un punto --- el volumen de la materia es de Pauli.
\end{solution}

\begin{solution}{exo:b3:quantum-statistics:12}
(a) Electrones activos: $g(E_{\text{F}})k_{\text{B}}T$; energía de cada uno:
$\sim k_{\text{B}}T$; derívese $E_{\text{extra}} \sim
g k_{\text{B}}^2T^2$. (b) $g(E_{\text{F}}) = 3N/2E_{\text{F}}$
da $C \sim 3Nk_{\text{B}}T/T_{\text{F}}$ --- el
$\pi^2/2$ exacto sustituye al 3. (c) Los espines volteables $\propto T$; cada uno
aporta una alineación $\propto \mu^2B/k_{\text{B}}T$: las $T$ se
cancelan --- la susceptibilidad de Pauli es plana allí donde la de Curie
(\cref{exo:b3:canonical-ensemble:11}) cae como $1/T$: un
diagnóstico de degeneración. (d) Los fermiones degenerados responden solo en
su superficie: multiplíquese cualquier respuesta clásica por $\sim
T/T_{\text{F}}$, o restrínjase a la capa de Fermi.
\end{solution}

\begin{solution}{pb:b3:quantum-statistics:1}
\textbf{1.} La energía liberada al ensamblar la estrella a partir de materia
dispersa --- negativa, y más profunda cuanto más compacta.
\textbf{2.} $E_{\text{F}} \sim (\hbar^2/m_{\text{e}})
(N_{\text{e}}/R^3)^{2/3}$.
\textbf{3.} $E_{\text{k}} \sim N_{\text{e}}E_{\text{F}} \sim
\hbar^2N_{\text{e}}^{5/3}/m_{\text{e}}R^2$.
\textbf{4.} Minimizar $A/R^2 - B/R$ da $R = 2A/B$: $R \sim
\hbar^2N_{\text{e}}^{-1/3}/(Gm_{\text{e}}m_{\text{n}}^2\mu_e^2)$.
\textbf{5.} $N_{\text{e}} \propto M$: $R \propto M^{-1/3}$ ---
añadir masa \emph{encoge} la estrella, de modo que la densidad y $E_{\text{F}}$
crecen sin límite: el sostén no relativista vive a
crédito.
\textbf{6.} $N_{\text{e}} = \num{6e56}$: $R \sim \qty{2000}{km}$
--- el orden correcto al lado de los \qty{5800}{km} de Sirio B (nuestro
escalado descartó factores de unas pocas unidades).
\textbf{7.} $E_{\text{F}} \sim m_{\text{e}}c^2$ en $n \sim
(m_{\text{e}}c/\hbar)^3 \approx \qty{2e37}{m^{-3}}$, es decir, $\rho
\sim \num{e10}$--$\num{e11}\,\unit{kg/m^3}$: el régimen de las
enanas más pesadas.
\textbf{8.} Cada electrón: $\epsilon \sim \hbar c\,n^{1/3}$;
en total: $\hbar cN_{\text{e}}^{4/3}/R$.
\textbf{9.} Con ambos términos $\propto 1/R$, el radio desaparece:
la estabilidad la decide si el coeficiente cinético supera
al gravitatorio --- un criterio de masa puro.
\textbf{10.} Iguálense: $N_{\text{e,crit}} \sim (\hbar c/
Gm_{\text{n}}^2\mu_e^2)^{3/2}$, y $M_{\text{Ch}} =
\mu_em_{\text{n}}N_{\text{e,crit}}$: la fórmula enunciada.
\textbf{11.} $\hbar c/Gm_{\text{n}}^2 = \num{1.7e38}$;
$M_{\text{Ch}} \sim (\num{1.7e38})^{3/2}m_{\text{n}}/4 \approx
\qty{9e29}{kg} \approx 0.5\,M_\odot$ --- el cálculo exacto
afina el orden de magnitud hasta $1.4\,M_\odot$.
\textbf{12.} $\hbar$ (la presión es cuántica), $c$ (su
ablandamiento relativista) y $G$ (el adversario): un número del tamaño de una
estrella construido con tres constantes microscópicas --- la mecánica cuántica
legislando para objetos de $10^{57}$ partículas.
\textbf{13.} La degeneración de neutrones, a un radio menor en
$\sim m_{\text{e}}/m_{\text{n}}$ (por la contabilidad de $\mu_e$):
kilómetros --- la estrella de neutrones.
\textbf{14.} Las mismas tres constantes con $m_{\text{n}}$
por doquier dan la misma escala de masa: las estrellas de neutrones se topan
con un techo cerca de $2\,M_\odot$; más allá, ninguna ley de presión gana ---
un agujero negro.
\textbf{15.} $GM^2/R \approx \qty{5e46}{J}$: varios cientos de veces
toda la producción del Sol en diez mil millones de años, liberada en segundos.
\textbf{16.} La supernova 1987A: dos docenas de neutrinos, llegados horas
antes de que la luz se avivara, portando justo la energía predicha
--- el colapso oído directamente.
\textbf{17.} Colapso de núcleo: el núcleo de hierro de una estrella masiva
implosiona más allá de todo suelo de presión; la envoltura sale despedida. Tipo
Ia: una \omterm{ex:b3:quantum-statistics:dwarf}{enana blanca} alimentada hasta el umbral enciende su carbono en un
fogonazo termonuclear que desliga la estrella entera.
\textbf{18.} La misma masa disparadora, el mismo combustible, la misma
energía liberada: curvas de luz casi idénticas --- bombas calibradas.
\textbf{19.} Conocido el brillo absoluto, el brillo observado
da la distancia por el inverso del cuadrado; sin
lo de ``estándar'', la escalera se viene abajo.
\textbf{20.} Demasiado débil significa demasiado lejos: la expansión se ha ido
acelerando --- energía oscura, apoyada probatoriamente en la
uniformidad de las velas (de ahí el cuidado con su física).
\textbf{21.} Principio de exclusión $\to$ \omterm{thm:b3:quantum-statistics:fermi}{presión de degeneración} $\to$ una
masa límite universal $\to$ explosiones estandarizadas $\to$ la
aceleración medida del universo.
\textbf{22.} Síganse las matemáticas adonde lleven: el
comportamiento ``absurdo'' de las estrellas más allá del límite resultó ser
el de las estrellas de neutrones y los agujeros negros --- ambos catalogados
hoy por millares.
\textbf{23.} En 1915 ninguna física conocida permitía a la materia una
tonelada por centímetro cúbico; en 1926, la estadística de Fermi--Dirac la
convirtió en el estado por defecto de cualquier gas frío y denso: la
observación esperó a la estadística, y no al revés.
\textbf{24.} Que las enanas existen de forma estable en todo un rango de masas
(la rama $R(M)$) y que la rama \emph{termina} en
$1.4\,M_\odot$: ambas cosas están dibujadas en ese histograma.
\textbf{25.} $(\hbar c/Gm_{\text{n}}^2)^{3/2}m_{\text{n}} \approx
1.4\,M_\odot$: tres constantes fijan la \omterm{ex:b3:quantum-statistics:dwarf}{enana blanca} más pesada ---
y con ella la masa que dispara las velas estándar que pesaron
el destino del universo.
\end{solution}
